\documentclass[a4paper,11pt]{article}
\usepackage{fontspec}
\usepackage[french]{babel}
\usepackage{amsmath,amssymb,amsthm}
\usepackage{geometry}
\usepackage{enumitem}
\usepackage{hyperref}
\usepackage{fancyhdr}
\usepackage{xcolor}
\usepackage{listings}
\usepackage{array}
\usepackage{booktabs}
\usepackage{multirow}
\usepackage{graphicx}

\geometry{margin=2cm}
\pagestyle{fancy}
\fancyhf{}
\fancyhead[C]{\textbf{STA211 — Analyse multi-blocs — Fiche récapitulative}}
\fancyfoot[C]{\thepage}
\renewcommand{\headrulewidth}{0.4pt}

\definecolor{codegreen}{rgb}{0.2,0.6,0.2}
\definecolor{codegray}{rgb}{0.5,0.5,0.5}
\definecolor{codepurple}{rgb}{0.58,0,0.82}
\definecolor{backcolour}{rgb}{0.95,0.95,0.92}

\lstdefinestyle{mystyle}{
    backgroundcolor=\color{backcolour},
    commentstyle=\color{codegreen},
    keywordstyle=\color{magenta},
    numberstyle=\tiny\color{codegray},
    stringstyle=\color{codepurple},
    basicstyle=\footnotesize\ttfamily,
    breakatwhitespace=false,
    breaklines=true,
    captionpos=b,
    keepspaces=true,
    numbers=left,
    numbersep=5pt,
    showspaces=false,
    showstringspaces=false,
    showtabs=false,
    tabsize=2
}
\lstset{style=mystyle}

\setlength{\parindent}{0pt}

\begin{document}

\begin{center}
{\LARGE\textbf{Analyse multi-blocs}} \\
\vspace{0.3cm}
STA211 — Apprentissage non supervisé — Fiche récapitulative \\
8 juin 2026
\end{center}

\vspace{0.3cm}
\rule{\textwidth}{0.4pt}
\vspace{0.3cm}

\tableofcontents
\newpage

% ============================================================
\section{Présentation}
% ============================================================

L'\textbf{analyse multi-blocs} regroupe un ensemble de méthodes factorielles permettant d'analyser simultanément plusieurs tableaux de données décrivant les mêmes individus (ou les mêmes variables). Chaque tableau est appelé un \textbf{bloc}.

\paragraph{Contexte typique :}
On dispose de $N$ tableaux $\mathbf{X}_1, \ldots, \mathbf{X}_N$, tous avec les mêmes $n$ individus en lignes, mais avec possiblement des variables différentes en colonnes ($p_1, \ldots, p_N$).

\paragraph{Exemple :}
Pour chaque département français, on dispose du nombre de crimes et délits dans 9 catégories, pour chaque année de 1974 à 1993. Chaque année est un \textbf{bloc} de 9 variables.

\paragraph{Objectifs :}
\begin{itemize}
    \item Comparer les groupes de variables entre eux
    \item Dégager une structure commune (compromis)
    \item Visualiser l'évolution des individus (trajectoires)
    \item Réduire la dimension
\end{itemize}

% ============================================================
\section{Phases de l'analyse multi-blocs}
% ============================================================

Toute analyse multi-blocs se décompose en \textbf{4 phases} :

\begin{enumerate}
    \item \textbf{Interstructure} : chaque tableau est résumé par un objet (centre de gravité, matrice de produits scalaires\ldots). On compare ces objets entre eux (ex. ACP sur les centres de gravité).
    \item \textbf{Compromis} : on recherche un espace commun qui résume l'ensemble des tableaux (maximisation de l'inertie projetée de tous les nuages).
    \item \textbf{Intrastructure} : on analyse le compromis (cercles des corrélations, contributions\ldots) pour décrire finement les ressemblances entre tableaux.
    \item \textbf{Trajectoires} : on projette les individus de chaque tableau dans le plan du compromis pour visualiser leur évolution.
\end{enumerate}

% ============================================================
\section{Méthodes}
% ============================================================

\subsection{Double ACP (DACP)}
\label{sec:dacp}

\paragraph{Quand l'utiliser ?}
Mêmes variables mesurées sur les mêmes individus à des instants différents ($p_t = p$ pour tout $t$, $n$ identique).

\paragraph{Principe :}
\begin{enumerate}
    \item \textbf{Interstructure} : ACP sur les centres de gravité $\bar{\mathbf{x}}_t$ de chaque tableau (dimension $N \times p$). Chaque année est un point dans $\mathbb{R}^p$.
    \item \textbf{Compromis} : on cherche les axes $\mathbf{v}_k$ maximisant la somme des inerties de projection de tous les nuages :
    \[
    \max_{\mathbf{v}_1,\ldots,\mathbf{v}_q} \sum_{t=1}^N \sum_{k=1}^q \mathbf{v}_k' \mathbf{V}_t \mathbf{v}_k
    \]
    où $\mathbf{V}_t$ est la matrice de variance-covariance du tableau $t$. Ceci revient à une ACP sur la matrice de variance-covariance \textbf{compromis} :
    \[
    \mathbf{V} = \sum_{t=1}^N \mathbf{V}_t
    \]
    \item \textbf{Intrastructure} : ACP sur les $N$ tableaux centrés concaténés verticalement (dimension $Nn \times p$). Les axes sont les mêmes que ceux du compromis.
    \item \textbf{Trajectoires} : projection des individus de chaque année dans le plan factoriel.
\end{enumerate}

\subsection{STATIS}
\paragraph{Quand l'utiliser ?}
Mêmes individus, variables différentes selon les blocs (ou l'inverse pour STATIS duale).

\paragraph{Principe :}
On ne peut plus comparer les centres de gravité (variables différentes). On utilise les \textbf{ matrices de produits scalaires} $\mathbf{W}_t = \mathbf{X}_t \mathbf{M}_t \mathbf{X}_t^\top$ de dimension $n \times n$, qui représentent les proximités entre individus pour chaque bloc.
\begin{enumerate}
    \item \textbf{Interstructure} : ACP sur l'espace des matrices $\mathbf{W}_t$ (produit scalaire de Hilbert-Schmidt).
    \item \textbf{Compromis} : combinaison linéaire optimale $\mathbf{W}_{\text{comp}} = \sum_{t=1}^N a_t \mathbf{W}_t$.
    \item \textbf{Intrastructure} : diagonalisation de $\mathbf{W}_{\text{comp}}$ (équivalente à une ACP).
\end{enumerate}

\subsection{AFM (Analyse Factorielle Multiple)}
\paragraph{Quand l'utiliser ?}
Mêmes individus, variables de types potentiellement différents (quantitatives, qualitatives). Elle équilibre le poids de chaque bloc en pondérant par la première valeur propre de l'ACP (ou ACM) de chaque bloc.

% ============================================================
\section{Exemple numérique : DACP sur les données de criminalité}
% ============================================================

\subsection{Données}
$n = 94$ départements français, $p = 9$ variables de criminalité mesurées chaque année de 1974 à 1993 ($N = 20$ tableaux).

Variables :
\begin{center}
\begin{tabular}{ll}
VO : vols et recels & FX : faux et escroqueries \\
DF : délits financiers & CH : chèques sans provision \\
CR : coups et règlements de comptes & DD : destructions et dégradations \\
ET : délits à la police des étrangers & ST : stupéfiants \\
\end{tabular}
\end{center}

\subsection{Phase 1 : Interstructure}

On calcule le centre de gravité de chaque année (vecteur de dimension 9), puis on effectue une ACP sur le tableau $20 \times 9$ centré.

\begin{center}
\begin{tabular}{lrr}
\hline
Axe & Valeur propre & \% inertie cumulé \\
\hline
1 & 6,83 & 75,9\% \\
2 & 1,69 & 94,7\% \\
3 & 0,24 & 97,4\% \\
\hline
\end{tabular}
\end{center}

Le premier plan explique \textbf{94,7\%} de l'inertie : les années sont très bien représentées.

\paragraph{Cercle des corrélations (interstructure) :}
\begin{itemize}
    \item Axe 1 : corrélé positivement avec VO, DD, ET, FX, CR, ST — \og effet taille \fg{} de la criminalité.
    \item Axe 2 : lié à CH et DF — oppose la criminalité financière au reste.
\end{itemize}

\paragraph{Plan des individus (années) :}
Les années s'alignent le long de l'axe 1, révélant une \textbf{évolution temporelle quasi linéaire} : la criminalité (VO, ST, ET, FX, CR, DD) a augmenté entre 1974 et 1993. Les années 1974--1978 sont stables, 1979--1988 montrent une forte progression, et 1989--1993 se stabilisent à un niveau élevé.

\subsection{Phase 2 : Compromis et Intrastructure}

On concatène les 20 tableaux centrés verticalement (matrice $20 \times 94 \times 9$, soit $1880 \times 9$) et on effectue une ACP.

\begin{center}
\begin{tabular}{lrr}
\hline
Axe & Valeur propre & \% inertie \\
\hline
1 & 4,18 & 32,2\% \\
2 & 2,36 & 18,2\% \\
3 & 1,93 & 14,9\% \\
\hline
\end{tabular}
\end{center}

\textbf{Interprétation :}
\begin{itemize}
    \item Axe 1 : oppose la criminalité contre les biens (VO, DD, CH) à la criminalité contre les personnes.
    \item Axe 2 : oppose criminalité financière (DF, FX) à la délinquance de voie publique.
\end{itemize}

\subsection{Phase 3 : Trajectoires}

On projette les individus (départements) de chaque année dans le plan 1--2 du compromis. On obtient des \textbf{trajectoires} qui montrent l'évolution de chaque département au cours du temps.

% ============================================================
\section{Code Python}
% ============================================================

\begin{lstlisting}[language=Python, caption=DACP avec scikit-learn]
import numpy as np
import pandas as pd
from sklearn.decomposition import PCA
from sklearn.preprocessing import StandardScaler

def dacp(list_of_dfs, n_components=2):
    """
    Double ACP sur une liste de DataFrames (meme variables, meme individus).
    Chaque DF = un bloc (ex. une annee).
    """
    n_blocs = len(list_of_dfs)
    p = list_of_dfs[0].shape[1]  # nombre de variables

    # --- Phase 1 : Interstructure ---
    # Centres de gravite de chaque bloc
    centers = np.array([df.mean(axis=0).values for df in list_of_dfs])
    centers_scaled = StandardScaler(with_std=False).fit_transform(centers)

    pca_inter = PCA()
    pca_inter.fit(centers_scaled)

    # --- Phase 2 : Compromis ---
    # Matrice de variance-covariance compromis
    V_sum = np.zeros((p, p))
    for df in list_of_dfs:
        df_centered = df - df.mean(axis=0)
        V_sum += df_centered.T @ df_centered / (len(df) - 1)
    V_sum /= n_blocs  # moyenne des matrices de covariance

    # ACP sur V_sum
    eigvals, eigvecs = np.linalg.eigh(V_sum)
    idx = np.argsort(eigvals)[::-1]
    eigvals = eigvals[idx]
    eigvecs = eigvecs[:, idx]

    # --- Phase 3 : Intrastructure ---
    # Concaténation verticale des tableaux centres
    all_centered = np.vstack([df - df.mean(axis=0) for df in list_of_dfs])
    scores = all_centered @ eigvecs[:, :n_components]

    # --- Phase 4 : Trajectoires ---
    n_individuals = list_of_dfs[0].shape[0]
    trajectories = np.zeros((n_individuals, n_blocs, n_components))
    for t, df in enumerate(list_of_dfs):
        centered = df - df.mean(axis=0)
        trajectories[:, t, :] = centered @ eigvecs[:, :n_components]

    return {
        'pca_inter': pca_inter,
        'eigvals': eigvals,
        'eigvecs': eigvecs,
        'scores': scores,
        'trajectories': trajectories
    }

# Exemple d'utilisation
if __name__ == "__main__":
    # Simulation de 3 blocs (annees) avec 10 individus et 4 variables
    np.random.seed(42)
    blocs = [pd.DataFrame(np.random.randn(10, 4) + t*0.1)
             for t in range(3)]
    res = dacp(blocs)
    print("Valeurs propres :", res['eigvals'][:3])
\end{lstlisting}

\newpage
\begin{lstlisting}[language=Python, caption=AFM avec prince]
# pip install prince

import pandas as pd
from prince import MFA

# Construction du jeu de donnees multi-blocs
# Simulons 2 blocs de variables
np.random.seed(123)
n = 50
bloc1 = pd.DataFrame({
    'V1': np.random.randn(n),
    'V2': np.random.randn(n),
    'V3': np.random.randn(n)
}, index=[f'I{i}' for i in range(n)])

bloc2 = pd.DataFrame({
    'V4': np.random.randn(n),
    'V5': np.random.randn(n)
}, index=[f'I{i}' for i in range(n)])

# Concaténation
X = pd.concat([bloc1, bloc2], axis=1)

# Groupes
groups = [3, 2]  # 3 var dans bloc1, 2 dans bloc2

# AFM
mfa = MFA(groups=groups, n_components=3)
mfa.fit(X)

# Résultats
print("Valeurs propres :", mfa.eigenvalues_)
print("Coordonnées des individus :")
print(mfa.row_coordinates(X).head())
\end{lstlisting}

% ============================================================
\section{Code R}
% ============================================================

\begin{lstlisting}[language=R, caption=DACP avec R]
# install.packages("FactoMineR")
library(FactoMineR)

# Simulation de 3 blocs (annees)
set.seed(42)
n <- 10  # individus
p <- 4   # variables
blocs <- list(
  data.frame(matrix(rnorm(n*p) + 0.1, n, p)),
  data.frame(matrix(rnorm(n*p) + 0.2, n, p)),
  data.frame(matrix(rnorm(n*p) + 0.3, n, p))
)

# Concaténation verticale pour DACP
X_all <- do.call(rbind, blocs)

# Variable indiquant le bloc
bloc_labels <- rep(1:3, each = n)

# DACP manuelle
centers <- t(sapply(blocs, colMeans))  # centres de gravite
pca_inter <- PCA(centers, graph = FALSE)
print(pca_inter$eig)

# Compromis
V_sum <- Reduce(`+`, lapply(blocs, cov))
eig <- eigen(V_sum)
print(eig$values[1:3])
\end{lstlisting}

\begin{lstlisting}[language=R, caption=AFM avec FactoMineR]
library(FactoMineR)

# Jeu de donnees : meme structure
n <- 50
bloc1 <- data.frame(
  V1 = rnorm(n), V2 = rnorm(n), V3 = rnorm(n)
)
bloc2 <- data.frame(
  V4 = rnorm(n), V5 = rnorm(n)
)
X <- cbind(bloc1, bloc2)

# AFM
res.afm <- MFA(X, group = c(3, 2),
               type = c("s", "s"),
               name.group = c("Bloc1", "Bloc2"),
               graph = FALSE)

# Valeurs propres
print(res.afm$eig)

# Cercle des correlations
plot(res.afm, choix = "var")

# Carte des individus
plot(res.afm, choix = "ind")

# Representations partielles
plot(res.afm, choix = "axes")
\end{lstlisting}

\newpage
% ============================================================
\section*{QCM — Vérifiez vos connaissances}
% ============================================================

\medskip
\textit{Pour chaque question, choisissez la bonne réponse. Les réponses sont en page \pageref{reponses-mb}.}

\begin{enumerate}

    \item Que signifie l'acronyme \textbf{DACP} ?
    \begin{enumerate}
        \item Double Analyse en Composantes Principales
        \item Discrimination par Analyse en Composantes Principales
        \item Décomposition en Axes Canoniques Principaux
        \item Analyse Dynamique en Composantes Principales
    \end{enumerate}

    \item Combien de phases comporte une analyse multi-blocs ?
    \begin{enumerate}
        \item 2
        \item 3
        \item 4
        \item 5
    \end{enumerate}

    \item La phase d'\textbf{interstructure} consiste à :
    \begin{enumerate}
        \item analyser la structure interne de chaque bloc séparément
        \item comparer les blocs entre eux via leurs centres de gravité ou matrices de produits scalaires
        \item calculer les trajectoires des individus
        \item effectuer une classification hiérarchique
    \end{enumerate}

    \item Dans la DACP, la matrice de variance-covariance \textbf{compromis} $\mathbf{V}$ s'obtient par :
    \begin{enumerate}
        \item $\mathbf{V} = \sum_{t=1}^N \mathbf{V}_t$
        \item $\mathbf{V} = \prod_{t=1}^N \mathbf{V}_t$
        \item $\mathbf{V} = \frac{1}{N}\sum_{t=1}^N \mathbf{V}_t$
        \item $\mathbf{V} = \sum_{t=1}^N \mathbf{V}_t^{-1}$
    \end{enumerate}

    \item La méthode \textbf{STATIS} utilise pour l'interstructure :
    \begin{enumerate}
        \item les centres de gravité de chaque tableau
        \item les matrices de produits scalaires $\mathbf{W}_t = \mathbf{X}_t \mathbf{M}_t \mathbf{X}_t^\top$
        \item les matrices de corrélation
        \item les valeurs propres de chaque tableau
    \end{enumerate}

    \item La DACP s'applique quand :
    \begin{enumerate}
        \item les individus sont différents d'un tableau à l'autre
        \item les variables sont différentes d'un tableau à l'autre
        \item les mêmes variables sont mesurées sur les mêmes individus
        \item les données sont qualitatives
    \end{enumerate}

    \item L'\textbf{AFM} (Analyse Factorielle Multiple) pondère chaque bloc par :
    \begin{enumerate}
        \item la somme de ses valeurs propres
        \item l'inverse de sa première valeur propre
        \item le nombre de variables du bloc
        \item le nombre d'individus
    \end{enumerate}

    \item Dans la DACP sur les données de criminalité, le premier axe de l'interstructure représente :
    \begin{enumerate}
        \item l'opposition entre crimes violents et délits financiers
        \item un effet taille de la criminalité qui augmente avec le temps
        \item la différence entre zones urbaines et rurales
        \item la saisonnalité des crimes
    \end{enumerate}

    \item La phase de \textbf{trajectoires} permet de :
    \begin{enumerate}
        \item visualiser l'évolution de chaque individu au fil des blocs
        \item tracer la frontière entre clusters
        \item estimer la qualité de la représentation
        \item déterminer le nombre optimal d'axes
    \end{enumerate}

    \item Quel package R permet de réaliser une AFM ?
    \begin{enumerate}
        \item ggplot2
        \item FactoMineR
        \item kohonen
        \item randomForest
    \end{enumerate}

\end{enumerate}

\newpage
\label{reponses-mb}

\begin{center}
{\Large\textbf{Réponses au QCM}}
\end{center}

\medskip

\begin{tabular}{cp{10cm}}
\toprule
\textbf{N°} & \textbf{Réponse} & \textbf{Explication} \\
\midrule
1 & \textbf{a} & DACP = Double Analyse en Composantes Principales. \\
2 & \textbf{c} & Les 4 phases sont : interstructure, compromis, intrastructure, trajectoires. \\
3 & \textbf{b} & L'interstructure compare les blocs (via centres de gravité pour DACP, matrices $\mathbf{W}_t$ pour STATIS). \\
4 & \textbf{a} & $\mathbf{V} = \sum_{t=1}^N \mathbf{V}_t$ (somme des matrices de variance-covariance de chaque bloc). \\
5 & \textbf{b} & STATIS utilise $\mathbf{W}_t = \mathbf{X}_t \mathbf{M}_t \mathbf{X}_t^\top$ (matrices $n \times n$). \\
6 & \textbf{c} & DACP : mêmes variables ($p_t = p$) sur les mêmes $n$ individus. \\
7 & \textbf{b} & L'AFM pondère chaque bloc par $1/\lambda_1^{(t)}$ (inverse de la 1\up{ère} valeur propre). \\
8 & \textbf{b} & Axe 1 = effet taille (corrélation avec VO, DD, ET, FX, CR, ST), croissance temporelle. \\
9 & \textbf{a} & Les trajectoires projettent les individus de chaque bloc dans le plan compromis. \\
10 & \textbf{b} & \texttt{FactoMineR} (fonction \texttt{MFA}) implémente l'AFM. \\
\bottomrule
\end{tabular}

\vspace{0.5cm}

\begin{center}
\rule{0.6\textwidth}{0.4pt}
\vspace{0.3cm}
\textit{Fin de la fiche.}
\end{center}

\end{document}
